\section{Local Proofs for the Singular Response Calculus}
\label{app:local_singular_response_proofs}

This appendix contains the local calculations used by the response theorem.
It has deliberately been reduced to eight load-bearing statements.  In
particular, it contains no independent constant records, no density-side
recovery assertion, and no high-frequency contour argument based only on a
polynomial resolvent bound.  Every estimate below states its input, output,
and dependence on the hypotheses in the main text.

\subsection{Impact, reflection, and the complete physical block}

\begin{proposition}[Local impact and reflection differentiation]
\label{prop:supp_local_impact_reflection_differentiation}
Let \(h_\theta\) be a regular one-step collision branch contained in a
homogeneous chart, and write the next impact as
\[
  G_\theta(z,\tau)
  =F_\theta(q(z)+\tau v(z))=0.
\]
If \(c_\theta=n_\theta\cdot v\ne0\) at the impact, then every mixed derivative
of the flight time of total order \(a\le J+2\) is a finite sum
\begin{equation}
\label{eq:supp_flight_time_derivative}
  \partial^a\tau_\theta
  =
  \sum_{\pi}
  c_\theta^{-1-|\pi|}
  P_{a,\pi}\bigl(
    \partial^{\le a}F_\theta,
    \partial^{\le a}q,
    \partial^{\le a}v,
    \partial^{<a}\tau_\theta
  \bigr),
\end{equation}
where the \(P_{a,\pi}\) are universal polynomials.  Derivatives of the
reflection law
\[
  R_\theta(q,v)
   =(q,v-2(v\cdot n_\theta(q))n_\theta(q))
\]
and of the branch Jacobian have the same form: a finite sum of products of
normal-frame, curvature, flight-time, and crossing-factor derivatives.  On
\(\Hstrip_k\) all coefficients are bounded by a polynomial in \(k\), with
degree depending only on \(a\) and the primitive \(C^{J+4}\) bounds.
\end{proposition}

\begin{proof}
Differentiate \(G_\theta(z,\tau_\theta(z))=0\).  At each order, isolate the
unique term
\(\partial_\tau G_\theta\,\partial^a\tau_\theta
  =c_\theta\,\partial^a\tau_\theta\);
all remaining terms contain lower flight-time derivatives.  This gives
\eqref{eq:supp_flight_time_derivative} by induction.  The reflection formula
is algebraic in \(v\) and \(n_\theta\), while the branch Jacobian is a finite
product of flight and reflection factors.  The homogeneous strip bounds
\(|c_\theta|^{-1}\le Ck^2\) and the finite derivative induction give the last
claim.  No summation over \(k\) is made in this proposition.
\end{proof}

\begin{theorem}[Verified Ward identities for complete invariant sources]
\label{thm:supp_complete_physical_ward_hierarchy}
Let \(\rho_s\) be the transported invariant collision density and let
\[
  G_j
  =
  \sum_{\ell=1}^j\binom j\ell
  (\partial_s^\ell\Ltransfer_s|_0)\rho^{(j-\ell)}_0 .
\]
Then, as a distributional identity on the smooth flux-seed/test core,
\begin{equation}
\label{eq:supp_global_ward_identity}
  G_j=(I-\Ltransfer_0)\rho^{(j)}_0,
  \qquad
  \langle G_j,1\rangle=0 .
\end{equation}
For the local cutoff model \(c_s(u,y)u^3H(u-\eta_s(y))\), the singular
endpoint terms vanish identically for \(j\le3\), whereas the leading
fourth-order endpoint is
\begin{equation}
\label{eq:supp_fourth_order_cubic_endpoint}
  -6\,\dot\eta_0^{\,4}c_0(0,y)\delta(u).
\end{equation}
Consequently the cubic flux zero alone proves no order-\(j\) endpoint
cancellation for \(j\ge4\).  Under the additional operator Ward identities
\eqref{eq:finite_physical_ward_conditions}, the complete physical order-\(j\)
coefficient satisfies
\[
  F_{j,m}(u,y;\rho,\psi)
  =u\,\widetilde F_{j,m}(u,y;\rho,\psi)
\]
and its dyadic block has the uniform tail
\begin{equation}
\label{eq:supp_complete_physical_block_tail}
  |\langle\mathcal T_{j,m}^{\rm phys}\rho,\psi\rangle|
  \le
  C_j2^{-m}
  \|\rho\|_{\mathfrak D_U^{(j+2)}}
  \|\psi\|_{\mathfrak A_U^{(j+3)}} .
\end{equation}
\end{theorem}

\begin{proof}
Differentiate
\(\Ltransfer_s\rho_s=\rho_s\) exactly \(j\) times and move the
\(\ell=0\) term to the left.  This gives
\eqref{eq:supp_global_ward_identity}; pairing with \(1\) gives zero because
\(\Ltransfer_s\) preserves total mass.

For the cutoff model, repeated differentiation gives sums of
\(u^3\delta^{(n)}(u)\) multiplied by parameter jets of \(\eta_s\).  The
distributional identity
\[
  u^3\delta^{(n)}(u)
  =
  \begin{cases}
    0,&n<3,\\
    -n(n-1)(n-2)\delta^{(n-3)}(u),&n\ge3
  \end{cases}
\]
proves the first three assertions and gives
\eqref{eq:supp_fourth_order_cubic_endpoint}.  Thus the remaining
operator-level cancellation is precisely the finite family in
\eqref{eq:finite_physical_ward_conditions}; it is not inferred from the
cubic zero.  If \(\dot\eta_0(y_0)\ne0\) and the localized complete coefficient
at that germ is nonzero, the fourth-order atom cannot vanish for every
localized seed/test.  Hence arbitrary-seed R3 is not a generic property of a
moving grazing face; only an exact assembled-source identity can remove this
obstruction.

Once these residual endpoint coefficients vanish,
\(F_j^{\rm reg}(0,y)=0\), and the fundamental theorem of calculus gives
\[
  F_j^{\rm reg}(u,y)
  =
  u\int_0^1\partial_uF_j^{\rm reg}(tu,y)\,dt .
\]
On the \(m\)-th blown-up dyadic block, \(u\asymp2^{-m}\).  The derivative and
holonomy bounds in (R3) therefore give
\eqref{eq:supp_complete_physical_block_tail}.  Artificial faces have already
cancelled pairwise before this estimate, so no unsigned strip moment is used.
\end{proof}

\subsection{One-wall and clean two-wall currents}

\begin{lemma}[Second derivative of one moving wall]
\label{lem:supp_second_order_single_wall_trace}
Let
\[
  g_s=g_0+s\dot g_0+\tfrac12s^2\ddot g_0+o(s^2)
\]
in a \(C^2\) chart, with \(Dg_0\ne0\).  Then
\begin{equation}
\label{eq:supp_second_heaviside_derivative}
  \partial_s^2H(g_s)|_0
  =
  \ddot g_0\,\delta(g_0)
  +\dot g_0^{\,2}\delta'(g_0).
\end{equation}
Equivalently, if
\(g_s=g_0-sh_1+\frac12s^2h_2+o(s^2)\), the coefficient of
\(\delta(g_0)\) is \(+h_2\).  In normal coordinates \(y=g_0\), its action on
a compactly supported amplitude \(a\) is
\[
  \langle\partial_s^2H(g_s)|_0,a\rangle
  =
  \int_{y=0}\ddot g_0\,a\,J\,d\sigma
  -
  \int_{y=0}\partial_y(\dot g_0^{\,2}aJ)\,d\sigma .
\]
\end{lemma}

\begin{proof}
Apply the one-variable distributional chain rule twice:
\(\partial_sH(g_s)=\dot g_s\delta(g_s)\).  A second derivative gives
\(\ddot g_s\delta(g_s)+\dot g_s^2\delta'(g_s)\).  The final formula is the
definition of \(\delta'\) after the coarea change of variables.
\end{proof}

\begin{lemma}[Repeated derivatives of one generator]
\label{lem:supp_repeated_generator_trace_derivatives}
Let \(g_s\) be a single \(C^J\) face generator with \(Dg_0\ne0\).  For
\(1\le m\le J\),
\begin{equation}
\label{eq:supp_repeated_face_faa_di_bruno}
  \partial_s^mH(g_s)|_0
  =
  \sum_{\ell=0}^{m-1}
  P_{m,\ell}(\dot g_0,\ldots,\partial_s^{m-\ell}g_0)
  \delta^{(\ell)}(g_0),
\end{equation}
where \(P_{m,\ell}\) is the corresponding Bell polynomial.  This is a normal
derivative of one conormal current, not a product
\(\delta(g_0)^q\).  For a smooth core density and test, its weak norm is
bounded by the order-\((m+2)\) flux-seed norm times the order-\((m+3)\)
one-sided trace-test norm.
\end{lemma}

\begin{proof}
Faà di Bruno's formula applied to the distribution \(H\) gives
\eqref{eq:supp_repeated_face_faa_di_bruno}.  In the normal coordinate
\(y=g_0\), move the \(\ell\) normal derivatives from
\(\delta^{(\ell)}(y)\) to the coefficient, density, test, and chart
Jacobian.  These are exactly the finite normal jets controlled by the stated
seed/test norms.  Completion follows after physical block assembly by
\cref{thm:supp_complete_physical_ward_hierarchy}.
\end{proof}

\begin{lemma}[Clean double-trace coarea formula]
\label{lem:supp_worked_double_trace_coarea}
Let \(g_\Gamma,g_\Lambda\) be two distinct \(C^J\) face generators in a
two-dimensional collision chart and suppose
\[
  \Delta_{\Gamma\Lambda}(z)
  =|\det(Dg_\Gamma(z),Dg_\Lambda(z))|
  \ge\delta>0
\]
on their common zero set.  Then for every compactly supported continuous
amplitude \(a\),
\begin{equation}
\label{eq:supp_clean_double_trace}
  \langle\delta(g_\Gamma)\delta(g_\Lambda),a\rangle
  =
  \sum_{z\in\{g_\Gamma=g_\Lambda=0\}}
  \frac{a(z)}{\Delta_{\Gamma\Lambda}(z)} .
\end{equation}
Parameter derivatives of this current are finite sums of derivatives of the
moving point, of \(a\), and of \(\Delta_{\Gamma\Lambda}^{-1}\).  They remain
point-current derivatives associated with the same labelled pair and do not
create a third delta factor.
\end{lemma}

\begin{proof}
The map \(z\mapsto(g_\Gamma(z),g_\Lambda(z))\) is a local diffeomorphism at
each common zero.  A partition of unity and the two-dimensional change of
variables formula give \eqref{eq:supp_clean_double_trace}.  The implicit
function theorem makes each labelled zero \(C^J\) in the parameter.
Differentiation of the finite point-evaluation formula proves the final
claim.
\end{proof}

\begin{proposition}[Determinant lower-bound alternatives]
\label{prop:supp_normal_determinant_lower_bound_cases}
Under (R2), every labelled distinct-face pair lies in exactly one of the
following cases.
\begin{enumerate}[label=\textup{(\roman*)},wide,labelindent=0pt]
  \item On the central region \(u\ge u_J\), the pair is admitted and
  \(\Delta_{\Gamma\Lambda}\ge c_J>0\).
  \item In the grazing collar \(0\le u<u_J\), distinct physical faces have no
  common zero; only one-face grazing currents are admitted.
  \item The pair has a boundary birth/death, is tangent, has a triple
  unfolding, or changes incidence.
  It is outside the clean response chart and contributes no two-wall atom to
  the theorem.
\end{enumerate}
There is no fourth case in which infinitely many determinant shells are
summed merely from a sublevel-width estimate.
\end{proposition}

\begin{proof}
The central bound and the empty grazing collar are the two quantitative parts
of (R2).  The simple normal-crossings hypothesis excludes triple-labelled
vertices and incidence changes.  A tangency or boundary birth is outside the
admitted atlas.  Finally, the example \(\Delta(t)=|t|\) has
\[
  \int_{2^{-p-1}}^{2^{-p}}\frac{dt}{\Delta(t)}=\log2
\]
for every \(p\); this proves that sublevel width alone cannot justify an
infinite-shell alternative.
The half-chart fold
\(g_1=y\), \(g_2=y-u^2+s\), \(u\ge0\), shows that a boundary determinant jet
\(\det(Dg_1,Dg_2)=2u\) would not suffice: an interior intersection is born for
\(s>0\).  This is why the empty-collar condition is explicit.
\end{proof}

\begin{proposition}[Finite-order clean two-wall current calculus]
\label{prop:supp_two_wall_higher_order_response_closure}
Assume (R2)--(R3), and let \(m\le J\).  The order-\(m\) derivative of a
resolved branch indicator on the simple two-dimensional atlas is a finite sum
of:
\begin{enumerate}[label=\textup{(\alph*)},wide,labelindent=0pt]
  \item branch-interior derivatives;
  \item the one-wall terms in
  \cref{lem:supp_repeated_generator_trace_derivatives};
  \item parameter derivatives of the clean point atoms in
  \cref{lem:supp_worked_double_trace_coarea}.
\end{enumerate}
The point atoms in (c) occur only in the central region \(u\ge u_J\); the
grazing collar has no distinct-face common zero.
For every finite decorated core \(\mathscr F\), after artificial faces are
cancelled and physical grazing terms are assembled,
\[
  \|\mathcal T_{D,\mathscr F}^{(m)}\rho\|_{\mathfrak X_U^{(m)}}
  \le C_{m,\mathscr F}\|\rho\|_{\mathfrak D_U^{(m+2)}} .
\]
The constant depends on the finite one-step atlas, the determinant lower
bound, the order-\(m\) face jets, and the Ward bound, but not on a determinant
shell cutoff.  A uniform countable-atlas completion is asserted only under
(S1).
\end{proposition}

\begin{proof}
Differentiate the pushforward of the resolved indicator current.  A point
with no active face gives (a), one active face gives (b), and two active faces
give (c).  Simple normal crossings exclude a fourth local incidence type.
The one-face and two-face formulas are the preceding lemmas.  Point atoms are
placed in the decorated ambient source scale of
\cref{def:ambient_conormal_source_scale}; they are not asserted to belong to
the ordinary Demers--Zhang weak stable-leaf completion.  Determinants are
controlled by
\cref{prop:supp_normal_determinant_lower_bound_cases}; their polynomial strip
loss is combined with all other physical terms before summation.
\Cref{thm:supp_complete_physical_ward_hierarchy} then supplies the dyadic
factor \(2^{-m}\), while adjacent artificial faces cancel with opposite
orientation.  On a finite core this proves the displayed estimate.  The
recursively weighted all-branch sum includes additional propagation weights;
its completion is precisely the independent mapping certificate (S1), and is
not inferred from the scalar \(2^{-m}\) tail.
\end{proof}

\subsection{Continuous-time boundary}

No differentiated flow theorem is part of this proof ledger.  The retained
fixed-table BDL input and the transverse/grazing/global-weight obstructions are
\cref{prop:fixed_table_bdl_correspondence,prop:flow_obstruction_trichotomy}.
A future proof must first construct the dynamical branch/preimage-weighted
common graph domain and its negative-order inverse-Laplace symbol estimate;
only then would a graph-domain Neumann argument and dyadic integration by
parts become applicable.
